\section*{致谢}
我们谨向以下诸位宝贵的意见表示感谢：
Sinya Aoki,
Masayuki Asakawa,
Michael G. Endres,
Kazuo Fujikawa,
Leonardo Giusti,
Shoji Hashimoto,
Tetsuo Hatsuda,
Etsuko Itou,
Yoshio Kikukawa,
Masakiyo Kitazawa,
Tetsuya Onogi,
Giancarlo Rossi,
Shoichi Sasaki,
Yusuke Taniguchi，
尤其感谢
Martin L\"uscher
还向我们提供了他的私人研究笔记。
H.~S. 的工作部分受科研经费（Grant-in-Aid for Scientific Research）23540330 资助。

\appendix

\section{$\text{MS}$ 方案中的单圈重正化}
\label{sec:A}
\subsection{参数与基本场}
规范耦合：
\begin{equation}
   g_0^2=\mu^{2\epsilon}g^2
   \left\{
   1+\frac{g^2}{(4\pi)^2}\left[-\frac{11}{3}C_2(G)+\frac{4}{3}T(R)N_{\mathrm{f}}\right]
   \frac{1}{\epsilon}+O(g^4)
   \right\}.
\label{eq:(A1)}
\end{equation}

费米子质量：
\begin{equation}
   m_0=m
   \left[
   1+\frac{g^2}{(4\pi)^2}C_2(R)(-3)
   \frac{1}{\epsilon}+O(g^4)
   \right].
\label{eq:(A2)}
\end{equation}

规范势（费曼规范）：
\begin{equation}
   A_\mu^a(x)=\left[1+\frac{g^2}{(4\pi)^2}C_2(G)(-1)
   \frac{1}{\epsilon}+O(g^4)\right]A_{\mu R}^a(x).
\label{eq:(A3)}
\end{equation}

费米子场：
\begin{equation}
   \psi(x)=
   \left[1+\frac{g^2}{(4\pi)^2}C_2(R)\left(-\frac{1}{2}\right)
   \frac{1}{\epsilon}+O(g^4)\right]
   \psi_R(x).
\label{eq:(A4)}
\end{equation}

\subsection{复合算符}
裸算符~\eqref{eq:(4.7)}--\eqref{eq:(4.11)} 与相应的重正化算符
\begin{align}
   \left\{\mathcal{O}_{1\mu\nu}\right\}_R(x)&\equiv
   \left\{F_{\mu\rho}^aF_{\nu\rho}^a\right\}_R(x),
\label{eq:(A5)}\\
   \left\{\mathcal{O}_{2\mu\nu}\right\}_R(x)&\equiv
   \delta_{\mu\nu}\left\{F_{\rho\sigma}^aF_{\rho\sigma}^a\right\}_R(x),
\label{eq:(A6)}\\
   \left\{\mathcal{O}_{3\mu\nu}\right\}_R(x)&\equiv
   \left\{\Bar{\psi}
   \left(\gamma_\mu\overleftrightarrow{D}_\nu
   +\gamma_\nu\overleftrightarrow{D}_\mu\right)
   \psi\right\}_R(x),
\label{eq:(A7)}\\
   \left\{\mathcal{O}_{4\mu\nu}\right\}_R(x)&\equiv
   \delta_{\mu\nu}\left\{\Bar{\psi}
   \overleftrightarrow{\Slash{D}}
   \psi\right\}_R(x),
\label{eq:(A8)}\\
   \left\{\mathcal{O}_{5\mu\nu}\right\}_R(x)&\equiv
   \delta_{\mu\nu}m\left\{\Bar{\psi}
   \psi\right\}_R(x)
\label{eq:(A9)}
\end{align}
通过
\begin{equation}
   \mathcal{O}_{i\mu\nu}(x)
   =Z_{ij}\left\{\mathcal{O}_{j\mu\nu}\right\}_R(x).
\label{eq:(A10)}
\end{equation}
相联系。$\mathcal{O}_{1\mu\nu}$ 与~$\mathcal{O}_{2\mu\nu}$
算符重正化中的胶子贡献已在文献~\cite{Suzuki:2013gza}中确定。进一步计算算符重正化的费米子
贡献（对应正文中的图~A18 与~A19），并计入规范耦合与波函数归一化［即
式~\eqref{eq:(A1)} 与~\eqref{eq:(A3)}］，我们得
\begin{align}
   Z_{11}
   &=1+\frac{g^2}{(4\pi)^2}C_2(G)
   \left(-\frac{11}{3}\right)
   \frac{1}{\epsilon}+O(g^4),
\label{eq:(A11)}\\
   Z_{12}
   &=\frac{g^2}{(4\pi)^2}C_2(G)\frac{11}{12}\frac{1}{\epsilon}+O(g^4),
\label{eq:(A12)}\\
   Z_{13}
   &=\frac{g^4}{(4\pi)^2}C_2(R)\frac{2}{3}\frac{1}{\epsilon}+O(g^6),
\label{eq:(A13)}\\
   Z_{14}
   &=\frac{g^4}{(4\pi)^2}C_2(R)\left(-\frac{1}{3}\right)\frac{1}{\epsilon}
   +O(g^6),
\label{eq:(A14)}\\
   Z_{15}
   &=\frac{g^4}{(4\pi)^2}C_2(R)(-3)\frac{1}{\epsilon}+O(g^6),
\label{eq:(A15)}
\end{align}
以及
\begin{align}
   Z_{21}
   &=0,
\label{eq:(A16)}\\
   Z_{22}
   &=1+O(g^4),
\label{eq:(A17)}\\
   Z_{23}
   &=0,
\label{eq:(A18)}\\
   Z_{24}
   &=O(g^6),
\label{eq:(A19)}\\
   Z_{25}
   &=\frac{g^4}{(4\pi)^2}C_2(R)(-12)\frac{1}{\epsilon}+O(g^6).
\label{eq:(A20)}
\end{align}
由这些关系，在单圈阶进一步有
\begin{align}
   &\delta_{\mu\nu}\left\{F_{\mu\rho}^aF_{\nu\rho}^a\right\}_R(x)
\notag\\
   &=\left[
   1+\frac{g^2}{(4\pi)^2}C_2(G)\left(\frac{11}{6}\right)\right]
   \left\{F_{\rho\sigma}^aF_{\rho\sigma}^a\right\}_R(x)
\notag\\
   &\qquad{}
   +\frac{g^4}{(4\pi)^2}C_2(R)\left(-\frac{2}{3}\right)
   \left\{\Bar{\psi}\overleftrightarrow{\Slash{D}}\psi\right\}_R(x)
   +\frac{g^4}{(4\pi)^2}C_2(R)(-6)
   m\left\{\Bar{\psi}\psi\right\}_R(x).
\label{eq:(A21)}
\end{align}

另一方面，图 B03、B04、B05、B06、B07 的计算结合波函数重正
化~\eqref{eq:(A4)} 给出
\begin{align}
   Z_{31}
   &=\frac{1}{(4\pi)^2}T(R)N_{\mathrm{f}}
   \frac{16}{3}
   \frac{1}{\epsilon}+O(g^2),
\label{eq:(A22)}\\
   Z_{32}
   &=\frac{1}{(4\pi)^2}T(R)N_{\mathrm{f}}
   \left(-\frac{4}{3}\right)
   \frac{1}{\epsilon}+O(g^2),
\label{eq:(A23)}\\
   Z_{33}
   &=1+\frac{g^2}{(4\pi)^2}C_2(R)\left(-\frac{8}{3}\right)\frac{1}{\epsilon}
   +O(g^4),
\label{eq:(A24)}\\
   Z_{34}
   &=\frac{g^2}{(4\pi)^2}C_2(R)\frac{4}{3}\frac{1}{\epsilon}
   +O(g^4),
\label{eq:(A25)}\\
   Z_{35}
   &=O(g^4),
\label{eq:(A26)}
\end{align}
以及
\begin{equation}
   Z_{41}=Z_{43}=0,\qquad Z_{42}=Z_{45}=O(g^4),\qquad Z_{44}=1+O(g^4).
\label{eq:(A27)}
\end{equation}
这些关系的自洽性给出
\begin{align}
   &\delta_{\mu\nu}\left\{\Bar{\psi}\gamma_\mu\overleftrightarrow{D}_\nu
   \psi\right\}_R(x)
\notag\\
   &=\left[1+\frac{g^2}{(4\pi)^2}C_2(R)\frac{4}{3}\right]
   \left\{\Bar{\psi}\overleftrightarrow{\Slash{D}}
   \psi\right\}_R(x)
   +\frac{1}{(4\pi)^2}T(R)N_{\mathrm{f}}\left(-\frac{4}{3}\right)
   \left\{F_{\rho\sigma}^aF_{\rho\sigma}^a\right\}_R(x).
\label{eq:(A28)}
\end{align}

最后，一个一般性的定理（例如参见文献~\cite{Collins:1984xc}）断言
\begin{equation}
   Z_{51}=Z_{52}=Z_{53}=Z_{54}=0,\qquad Z_{55}=1.
\label{eq:(A29)}
\end{equation}

\section{积分公式}
\label{sec:B}
\begin{equation}
   \int_\ell\frac{1}{(\ell^2)^\alpha}\mathrm{e}^{-s\ell^2}
   =\frac{\Gamma(D/2-\alpha)}{(4\pi)^{D/2}\Gamma(D/2)}s^{\alpha-D/2}.
\label{eq:(B1)}
\end{equation}

\begin{equation}
   \int_k\int_\ell\frac{\mathrm{e}^{-sk^2-u\ell^2-v(k+\ell)^2}}{k^2}
   =\frac{1}{(4\pi)^D(D/2-1)(u+v)}(su+uv+vs)^{1-D/2}.
\label{eq:(B2)}
\end{equation}

\begin{align}
   &\int_\ell \mathrm{e}^{-s\ell^2}=\frac{1}{(4\pi)^{D/2}}s^{-D/2}.
\label{eq:(B3)}\\
   &\int_\ell \mathrm{e}^{-s\ell^2}\ell_\mu\ell_\nu
   =\frac{1}{(4\pi)^{D/2}}s^{-D/2-1}\frac{1}{2}\delta_{\mu\nu}.
\label{eq:(B4)}\\
   &\int_\ell \mathrm{e}^{-s\ell^2}
   \ell_\mu\ell_\nu\ell_\rho\ell_\sigma
   =\frac{1}{(4\pi)^{D/2}}s^{-D/2-2}
   \frac{1}{4}\left(
   \delta_{\mu\nu}\delta_{\rho\sigma}
   +\delta_{\mu\rho}\delta_{\nu\sigma}
   +\delta_{\mu\sigma}\delta_{\nu\rho}
   \right).
\label{eq:(B5)}\\
   &\int_\ell \mathrm{e}^{-s\ell^2}
   \ell_\mu\ell_\nu\ell_\rho\ell_\sigma\ell_\alpha\ell_\beta
   =\frac{1}{(4\pi)^{D/2}}s^{-D/2-3}\frac{1}{8}
   \left(
   \delta_{\mu\nu}\delta_{\rho\sigma}\delta_{\alpha\beta}
   +\text{$14$ 个置换}\right).
\label{eq:(B6)}
\end{align}

\begin{align}
   &\int_\ell\frac{1}{\ell^2}\,
   \mathrm{e}^{-s\ell^2}=\frac{1}{(4\pi)^{D/2}}\frac{1}{D/2-1}s^{-D/2+1}.
\label{eq:(B7)}\\
   &\int_\ell\frac{1}{\ell^2}\,\mathrm{e}^{-s\ell^2}\ell_\mu\ell_\nu
   =\frac{1}{(4\pi)^{D/2}}\frac{1}{D}s^{-D/2}\delta_{\mu\nu}.
\label{eq:(B8)}\\
   &\int_\ell\frac{1}{\ell^2}\,\mathrm{e}^{-s\ell^2}
   \ell_\mu\ell_\nu\ell_\rho\ell_\sigma
   =\frac{1}{(4\pi)^{D/2}}\frac{1}{2(D+2)}s^{-D/2-1}
   \left(
   \delta_{\mu\nu}\delta_{\rho\sigma}
   +\delta_{\mu\rho}\delta_{\nu\sigma}
   +\delta_{\mu\sigma}\delta_{\nu\rho}
   \right).
\label{eq:(B9)}\\
   &\int_\ell\frac{1}{\ell^2}\,\mathrm{e}^{-s\ell^2}
   \ell_\mu\ell_\nu\ell_\rho\ell_\sigma\ell_\alpha\ell_\beta
\notag\\
   &=\frac{1}{(4\pi)^{D/2}}\frac{1}{4(D+4)}s^{-D/2-2}
   \left(
   \delta_{\mu\nu}\delta_{\rho\sigma}\delta_{\alpha\beta}
   +\text{$14$ 个置换}\right).
\label{eq:(B10)}
\end{align}

\section{$\zeta_{1j}^{(1)}(t)$ 的胶子贡献}
\label{sec:C}
$\zeta_{1j}^{(1)}(t)$ 的胶子贡献列于表~\ref{table:0}。
\begin{table}
\caption{$\zeta_{1j}^{(1)}$，单位为 $\frac{g_0^2}{(4\pi)^2}C_2(G)$。
带~$*$ 的数值系相对本论文先前版本的更正。
}
\label{table:0}
\begin{center}
\renewcommand{\arraystretch}{2.2}
\setlength{\tabcolsep}{20pt}
\begin{tabular}{crr}
\toprule
 diagram & \multicolumn{1}{c}{$\zeta_{11}^{(1)}(t)$}
 & \multicolumn{1}{c}{$\zeta_{12}^{(1)}(t)$} \\
\midrule
 A03  & $0$ & $0$ \\
 A04  & $-3\epsilon(t)^{-1}-1$ & $0$ \\
 A05  & $-\dfrac{7}{36}$ & $-\dfrac{49}{144}$ \\
 A06  & $2\epsilon(t)^{-1}-\dfrac{1}{2}$ & $0$ \\
 A07  & $\dfrac{19}{288}$ & $\dfrac{121}{384}$ \\
 %A08  & $\dfrac{47}{96}$ & $\dfrac{53}{128}$ \\
 A08  & $\dfrac{35}{96}^*$ & $\dfrac{143}{384}^*$ \\
 A09  & $-\dfrac{25}{8}$ & $0$ \\
 A11  & $\dfrac{1}{3}\epsilon(t)^{-1}-\dfrac{17}{36}$ & $\dfrac{7}{12}\epsilon(t)^{-1}+\dfrac{1}{144}$ \\
 A13  & $-\dfrac{5}{3}\epsilon(t)^{-1}+\dfrac{25}{36}$ & $-\dfrac{3}{2}\epsilon(t)^{-1}-\dfrac{29}{16}$ \\
 A14  & $3\epsilon(t)^{-1}+3$ & $0$ \\
 A15  & $3\epsilon(t)^{-1}+\dfrac{5}{2}$ & $0$ \\
 %A16  & $\dfrac{7}{4}$ & $\dfrac{31}{24}$ \\
 A16  & $1^*$ & $\dfrac{31}{24}$ \\
\bottomrule
\end{tabular}
\end{center}
\end{table}

\section{我们计算方案的合理性论证}
\label{sec:D}
本附录详细解释：为何式~\eqref{eq:(4.25)} 中的单圈匹配系数~$\zeta^{(1)}_{ij}(t)$
可以在不计算关联函数、即把式~\eqref{eq:(4.26)} 和~\eqref{eq:(4.29)}
中的算符~$\Tilde{\mathcal{O}}_{i\mu\nu}(t,x)$ 换成式%
~\eqref{eq:(4.6)}--\eqref{eq:(4.10)} 中相应裸算符后的
\begin{equation}
   \left\langle\mathcal{O}_{i\mu\nu}(x)
   A_\beta^b(y)A_\gamma^c(z)\right\rangle,
\label{eq:(D1)}
\end{equation}
以及
\begin{equation}
   \left\langle\mathcal{O}_{i\mu\nu}(x)
   \psi(y)\Bar{\psi}(z)\right\rangle.
\label{eq:(D2)}
\end{equation}
的情形下加以确定。我们要论证：只要遵循正文中所采用的针对 IR 发散的正规化方案，就可以把这
些关联函数的贡献完全略去。

式~\eqref{eq:(D1)} 与~\eqref{eq:(D2)} 可以类似地处理，故考虑式~\eqref{eq:(D1)}。
任取一个对~\eqref{eq:(D1)} 有贡献的单圈 1PI 费曼图，例如图~\ref{fig:22}
中的图~A11。这是 $D$ 维规范理论中的图，因而只含普通（即实心圆）顶点与普通（即在
式~\eqref{eq:(3.17)} 中无高斯因子）的传播子。把该图对~\eqref{eq:(D1)}
的贡献写成式~\eqref{eq:(4.27)}，则顶点部
分~$\mathcal{M}_{\mu\nu,\beta\gamma}(k,\ell)$ 由单圈积分给出：
\begin{equation}
   \mathcal{M}_{\mu\nu,\beta\gamma}(k,\ell)=\int_p\,
   \mathcal{I}_{\mu\nu,\beta\gamma}(p;k,\ell),
\label{eq:(D3)}
\end{equation}
其中量纲分析表明被积函数~$\mathcal{I}_{\mu\nu,\beta\gamma}(p;k,\ell)$ 的质量量纲为%
~$-2$。一般而言该积分 UV 发散，而且若进一步对外动量 $k$ 与~$\ell$ 作泰勒展开，还会表现出
IR 发散。

其次注意，对 $D$ 维规范理论中的上述 1PI 图，总存在结构完全相同（即此处图~\ref{fig:22}
中的图~A11）的对应流费曼图。把其对式~\eqref{eq:(4.27)} 顶点部
分~$\mathcal{M}_{\mu\nu,\beta\gamma}(k,l)$ 的贡献写成
\begin{equation}
   \mathcal{M}_{\mu\nu,\beta\gamma}(k,\ell)=\int_p\,
   \mathcal{I}_{\mu\nu,\beta\gamma}(p;k,\ell;t),
\label{eq:(D4)}
\end{equation}
其中顶点的费曼规则与式~\eqref{eq:(D3)} 的完全相同，而传播子除式~\eqref{eq:(3.18)}
中依赖于插入复合算符~$\Tilde{\mathcal{O}}_{i\mu\nu}(t,x)$ 流时间的高斯因子外也完全相同；
我们已显式标出被积函数对~$t$ 的依赖。

现在，在\emph{单圈层次\/}上可以看到，与~$\zeta^{(1)}_{ij}(t)$ 相关的是式%
~\eqref{eq:(D4)} 与~\eqref{eq:(D3)} 之差：
\begin{equation}
   \int_p\,\left[
   \mathcal{I}_{\mu\nu,\beta\gamma}(p;k,\ell;t)
   -\mathcal{I}_{\mu\nu,\beta\gamma}(p;k,\ell)\right].
\label{eq:(D5)}
\end{equation}
由上述两个被积函数的结构，（对数型）IR 发散在该差中被消去，于是我们可以对外动量~$k$
和~$\ell$ 作泰勒展开。$O(k,\ell)$ 项的系数由
\begin{equation}
   C\int_p\frac{\mathrm{e}^{-2tp^2}-1}{(p^2)^2}
   =C\frac{1}{(4\pi)^{D/2}}\frac{4}{(D-2)(D-4)}(2t)^{2-D/2},
\label{eq:(D6)}
\end{equation}
给出，其中 $C$ 是常数，并引入了复维度 $D$ \emph{以正规化 UV 发散}；最后的表达式是由复区域
$2<\re(D)<4$ 解析延拓得到的。组合~\eqref{eq:(D6)}（取 $D\to4$）
正是所论图对 $\zeta_{ij}^{(1)}$ 贡献的相关部分。

现在，尽管上述计算是正规的，却存在一种以少得多的工夫重现式~\eqref{eq:(D6)}
的“便捷方法”；论证如下。

只考虑流费曼图~\eqref{eq:(D4)} 的贡献，并把被积函数对外动量 $k$
和~$\ell$ 作泰勒展开。这样便会遇到 IR 发散。然后引入复维度 $D$ \emph{以正规化 IR 发散}。
注意这是可行的，因为积分~\eqref{eq:(D4)} 得益于高斯因子而 UV 有限；总存在使积分良定义的
$D$ 的复区域。此时 $O(k,\ell)$ 项的系数为
\begin{equation}
   C\int_p\frac{\mathrm{e}^{-2tp^2}}{(p^2)^2}
   =C\frac{1}{(4\pi)^{D/2}}\frac{4}{(D-2)(D-4)}(2t)^{2-D/2}.
\label{eq:(D7)}
\end{equation}
这是由复区域~$4<\re(D)$ 解析延拓得到的表达式；然而它与式~\eqref{eq:(D6)}
完全相同。

于是我们得到了如下便捷方法：完全不去计算带裸算符的关联函数，即
式~\eqref{eq:(D1)} 与~\eqref{eq:(D2)}；对带流动算符的关联函数，即
式~\eqref{eq:(4.26)} 与~\eqref{eq:(4.29)}，只须简单地把费曼积分的被积函数对外动量（及费米
子质量）作泰勒展开；由此产生的 IR 发散用“维数正规化”、即自 $\re(D)>4$
的解析延拓来正规化。这样得到的单圈匹配系数~$\zeta_{ij}^{(1)}$
的结果，与正规地把式~\eqref{eq:(D1)} 和~\eqref{eq:(D2)} 的贡献计入的计算所得完全相同。
这一便捷方法正是我们在正文中所使用的计算方法。

\section{修改后的能动量张量}
\label{sec:E}
第~\ref{sec:2} 节的论证表明能动量张量~\eqref{eq:(2.7)} 满足 Ward--Takahashi 关系
\begin{equation}
   \left\langle\partial_\mu T_{\mu\nu}(x)A_\rho(y)\dotsm\right\rangle
   =-\delta(x-y)
   \left\langle\left[\partial_\nu A_\rho(y)
   -D_\rho A_\nu(y)\right]\dotsm\right\rangle+\dotsb,
\label{eq:(E1)}
\end{equation}
以及
\begin{equation}
   \left\langle\partial_\mu T_{\mu\nu}(x)\psi(y)\dotsm\right\rangle
   =-\delta(x-y)
   \left\langle D_\nu \psi(y)\dotsm\right\rangle+\dotsb.
\label{eq:(E2)}
\end{equation}
特别地，右边不存在正比于~$\partial_\nu\delta(x-y)$ 的项。这些关系导致
\begin{align}
   &\left\langle\partial_\mu\left[x_\nu T_{\mu\nu}(x)\right]
   A_\rho(y)\dotsm\right\rangle
\notag\\
   &=\left\langle T_{\mu\mu}(x)
   A_\rho(y)\dotsm\right\rangle
   -x_\nu\delta(x-y)
   \left\langle\left[\partial_\nu A_\rho(y)
   -D_\rho A_\nu(y)\right]\dotsm\right\rangle+\dotsb,
\label{eq:(E3)}
\end{align}
以及
\begin{equation}
   \left\langle\partial_\mu\left[x_\nu T_{\mu\nu}(x)\right]
   \psi(y)\dotsm\right\rangle
   =\left\langle T_{\mu\mu}(x)
   \psi(y)\dotsm\right\rangle
   -x_\nu\delta(x-y)
   \left\langle D_\nu \psi(y)\dotsm\right\rangle+\dotsb.
\label{eq:(E4)}
\end{equation}
对全空间积分这些关系，我们得到
\begin{align}
   \int\mathrm{d}^Dx\,
   \left\langle T_{\mu\mu}(x)
   A_\rho(y)\dotsm\right\rangle
   &=\left\langle\left[y_\nu\partial_\nu A_\rho(y)
   -y_\nu D_\rho A_\nu(y)\right]\dotsm\right\rangle+\dotsb
\notag\\
   &=\left\langle\left\{\left(y_\nu\partial_\nu+1\right) A_\rho(y)
   -D_\rho\left[y_\nu A_\nu(y)\right]\right\}\dotsm\right\rangle+\dotsb,
\label{eq:(E5)}
\end{align}
以及
\begin{align}
   \int\mathrm{d}^Dx\,\left\langle T_{\mu\mu}(x)
   \psi(y)\dotsm\right\rangle
   &=\left\langle y_\nu D_\nu \psi(y)\dotsm\right\rangle+\dotsb
\notag\\
   &=\left\langle\left[y_\nu\partial_\nu\psi(y)+y_\nu A_\nu(y)\psi(y)\right]
   \dotsm\right\rangle+\dotsb.
\label{eq:(E6)}
\end{align}

现在我们看到，式~\eqref{eq:(E5)} 与朴素预期的关系
\begin{equation}
   \int\mathrm{d}^Dx\,\left\langle T_{\mu\mu}(x)
   \mathcal{O}(y)\dotsm\right\rangle
   =\left\langle\left(y_\nu\partial_\nu+d_\mathcal{O}\right)
   \mathcal{O}(y)\dotsm\right\rangle+\dotsb,
\label{eq:(E7)}
\end{equation}
相符——其中 $d_{\mathcal{O}}$ 是场~$\mathcal{O}$ 的标度量纲——差别只在规范变换之内；
但式~\eqref{eq:(E6)} 并非如此。人们反而会预期
\begin{equation}
   \int\mathrm{d}^Dx\,\left\langle\Tilde{T}_{\mu\mu}(x)
   \psi(y)\dotsm\right\rangle
   =\left\langle\left(y_\nu\partial_\nu+\frac{D-1}{2}\right)\psi(y)
   \dotsm\right\rangle+\dotsb,
\label{eq:(E8)}
\end{equation}
（至多差一规范变换）。利用可以把运动方程加进能动量张量的自由，可以构造这样一个修改后的能
动量张量。可以看到，若
\begin{equation}
   \left\langle\partial_\mu\Tilde{T}_{\mu\nu}(x)\psi(y)\dotsm\right\rangle
   =-\delta(x-y)
   \left\langle D_\nu \psi(y)\dotsm\right\rangle
   +\frac{D-1}{2D}\partial_\nu^x\delta(x-y)
   \left\langle\psi(y)\dotsm\right\rangle+\dotsb.
\label{eq:(E9)}
\end{equation}
则式~\eqref{eq:(E8)} 得以实现。可以看出最后一项不影响与平移及转动（Poincar\'e
变换）相联系的 Ward--Takahashi 关系。这一修改由
\begin{equation}
   \Tilde{T}_{\mu\nu}(x)=T_{\mu\nu}(x)
   +\frac{D-1}{2D}
   \delta_{\mu\nu}
   \Bar{\psi}(x)\left(\overleftrightarrow{\Slash{D}}+2m_0\right)\psi(x).
\label{eq:(E10)}
\end{equation}
完成。这个满足朴素预期关系~\eqref{eq:(E7)} 的修改能动量张量使迹反常的推导更为简洁。
用流动算符的语言，此修改相当于给组合~\eqref{eq:(5.1)} 乘上
\begin{equation}
   \frac{D-1}{2D}\to\frac{3}{8}.
\end{equation}
然而我们注意到：当能动量张量在位置空间中与其他算符分离时，该修改不产生任何影响，因为它正
比于运动方程。
\begin{thebibliography}{00}

%\cite{Callan:1970ze}
\bibitem{Callan:1970ze} 
  C.~G.~Callan, Jr., S.~R.~Coleman, and R.~Jackiw,
  %``A New improved energy - momentum tensor,''
  Ann.\ Phys.\  {\bf 59}, 42 (1970).
  %%CITATION = APNYA,59,42;%%
  %664 citations counted in INSPIRE as of 07 Mar 2014

%\cite{Coleman:1970je}
\bibitem{Coleman:1970je} 
  S.~R.~Coleman and R.~Jackiw,
  %``Why dilatation generators do not generate dilatations?,''
  Ann.\ Phys.\  {\bf 67}, 552 (1971).
  %%CITATION = APNYA,67,552;%%
  %263 citations counted in INSPIRE as of 07 Mar 2014

%\cite{Freedman:1974gs}
\bibitem{Freedman:1974gs} 
  D.~Z.~Freedman, I.~J.~Muzinich and E.~J.~Weinberg,
  %``On the Energy-Momentum Tensor in Gauge Field Theories,''
  Ann.\ Phys.\  {\bf 87}, 95 (1974).
  %%CITATION = APNYA,87,95;%%
  %79 citations counted in INSPIRE as of 07 Mar 2014

%\cite{Joglekar:1975jm}
\bibitem{Joglekar:1975jm} 
  S.~D.~Joglekar,
  %``An Energy-Momentum Tensor in Gauge Theories,''
  Ann.\ Phys.\  {\bf 100}, 395 (1976);
  {\bf 102}, 594 (1976) [erratum].
  %%CITATION = APNYA,100,395;%%
  %10 citations counted in INSPIRE as of 07 Mar 2014

%\cite{Caracciolo:1988hc}
\bibitem{Caracciolo:1988hc} 
  S.~Caracciolo, G.~Curci, P.~Menotti and A.~Pelissetto,
  %``The Energy Momentum Tensor on the Lattice: The Scalar Case,''
  Nucl.\ Phys.\ B {\bf 309}, 612 (1988).
  %%CITATION = NUPHA,B309,612;%%
  %15 citations counted in INSPIRE as of 07 Mar 2014

%\cite{Caracciolo:1989pt}
\bibitem{Caracciolo:1989pt} 
  S.~Caracciolo, G.~Curci, P.~Menotti and A.~Pelissetto,
  %``The Energy Momentum Tensor for Lattice Gauge Theories,''
  Ann.\ Phys.\  {\bf 197}, 119 (1990).
  %%CITATION = APNYA,197,119;%%
  %24 citations counted in INSPIRE as of 07 Mar 2014

%\cite{Suzuki:2013gza}
\bibitem{Suzuki:2013gza} 
  H.~Suzuki,
  %``Energy–momentum tensor from the Yang–Mills gradient flow,''
  PTEP {\bf 2013}, 083B03 (2013)
  [arXiv:1304.0533 [hep-lat]].
  %%CITATION = ARXIV:1304.0533;%%
  %5 citations counted in INSPIRE as of 07 Mar 2014

%\cite{Luscher:2010iy}
\bibitem{Luscher:2010iy} 
  M.~L\"uscher,
  %``Properties and uses of the Wilson flow in lattice QCD,''
  JHEP {\bf 1008}, 071 (2010)
  [arXiv:1006.4518 [hep-lat]].
  %%CITATION = ARXIV:1006.4518;%%
  %67 citations counted in INSPIRE as of 07 Mar 2014

%\cite{Luscher:2011bx}
\bibitem{Luscher:2011bx} 
  M.~L\"uscher and P.~Weisz,
  %``Perturbative analysis of the gradient flow in non-abelian gauge theories,''
  JHEP {\bf 1102}, 051 (2011)
  [arXiv:1101.0963 [hep-th]].
  %%CITATION = ARXIV:1101.0963;%%
  %27 citations counted in INSPIRE as of 07 Mar 2014

%\cite{Luscher:2013cpa}
\bibitem{Luscher:2013cpa} 
  M.~L\"uscher,
  %``Chiral symmetry and the Yang--Mills gradient flow,''
  JHEP {\bf 1304}, 123 (2013)
  [arXiv:1302.5246 [hep-lat]].
  %%CITATION = ARXIV:1302.5246;%%
  %11 citations counted in INSPIRE as of 07 Mar 2014

%\cite{Giusti:2010bb}
\bibitem{Giusti:2010bb} 
  L.~Giusti and H.~B.~Meyer,
  %``Thermal momentum distribution from path integrals with shifted boundary conditions,''
  Phys.\ Rev.\ Lett.\  {\bf 106}, 131601 (2011)
  [arXiv:1011.2727 [hep-lat]].
  %%CITATION = ARXIV:1011.2727;%%
  %13 citations counted in INSPIRE as of 07 Mar 2014

%\cite{Giusti:2012yj}
\bibitem{Giusti:2012yj} 
  L.~Giusti and H.~B.~Meyer,
  %``Implications of Poincare symmetry for thermal field theories in finite-volume,''
  JHEP {\bf 1301}, 140 (2013)
  [arXiv:1211.6669 [hep-lat]].
  %%CITATION = ARXIV:1211.6669;%%
  %9 citations counted in INSPIRE as of 07 Mar 2014

%\cite{Robaina:2013zmb}
\bibitem{Robaina:2013zmb} 
  D.~Robaina and H.~B.~Meyer,
  %``Renormalization of the momentum density on the lattice using shifted boundary conditions,''
  PoS LATTICE {\bf 2013}, 323 (2014)
  [arXiv:1310.6075 [hep-lat]].
  %%CITATION = ARXIV:1310.6075;%%
  %8 citations counted in INSPIRE as of 15 juin 2015

%\cite{Giusti:2013sqa}
\bibitem{Giusti:2013sqa} 
  L.~Giusti and H.~B.~Meyer,
  %``Thermal field theories and shifted boundary conditions,''
  PoS LATTICE {\bf 2013}, 214 (2013)
  [arXiv:1310.7818 [hep-lat]].
  %%CITATION = ARXIV:1310.7818;%%
  %1 citations counted in INSPIRE as of 07 Mar 2014

%\cite{Giusti:2013mxa}
\bibitem{Giusti:2013mxa} 
  L.~Giusti and M.~Pepe,
  %``Measuring the entropy from shifted boundary conditions,''
  arXiv:1311.1012 [hep-lat].
  %%CITATION = ARXIV:1311.1012;%%
  %4 citations counted in INSPIRE as of 15 Jun 2015

%\cite{Giusti:2014ila}
\bibitem{Giusti:2014ila} 
  L.~Giusti and M.~Pepe,
  %``Equation of state of a relativistic theory from a moving frame,''
  Phys.\ Rev.\ Lett.\  {\bf 113}, 031601 (2014)
  [arXiv:1403.0360 [hep-lat]].
  %%CITATION = ARXIV:1403.0360;%%
  %8 citations counted in INSPIRE as of 15 juin 2015

%\cite{Suzuki:2012wx}
\bibitem{Suzuki:2012wx} 
  H.~Suzuki,
  %``Remark on the energy-momentum tensor in the lattice formulation of 4D $\mathcal{N}=1$ SYM,''
  Phys.\ Lett.\ B {\bf 719}, 435 (2013)
  [arXiv:1209.5155 [hep-lat]].
  %%CITATION = ARXIV:1209.5155;%%
  %4 citations counted in INSPIRE as of 07 Mar 2014

%\cite{Luscher:2013vga}
\bibitem{Luscher:2013vga} 
  M.~L\"uscher,
  %``Future applications of the Yang-Mills gradient flow in lattice QCD,''
  PoS LATTICE {\bf 2013}, 016 (2014)
  [arXiv:1308.5598 [hep-lat]].
  %%CITATION = ARXIV:1308.5598;%%
  %28 citations counted in INSPIRE as of 15 juin 2015

%\cite{Borsanyi:2012zs}
\bibitem{Borsanyi:2012zs} 
  S.~Bors\'anyi, S.~D\"urr, Z.~Fodor, C.~Hoelbling, S.~D.~Katz, S.~Krieg, T.~Kurth and L.~Lellouch {\it et al.},
  %``High-precision scale setting in lattice QCD,''
  JHEP {\bf 1209}, 010 (2012)
  [arXiv:1203.4469 [hep-lat]].
  %%CITATION = ARXIV:1203.4469;%%
  %88 citations counted in INSPIRE as of 15 Jun 2015

%\cite{Borsanyi:2012zr}
\bibitem{Borsanyi:2012zr} 
  S.~Bors\'anyi, S.~D\"urr, Z.~Fodor, S.~D.~Katz, S.~Krieg, T.~Kurth, S.~Mages and A.~Sch\"afer {\it et al.},
  %``Anisotropy tuning with the Wilson flow,''
  arXiv:1205.0781 [hep-lat].
  %%CITATION = ARXIV:1205.0781;%%
  %10 citations counted in INSPIRE as of 15 juin 2015

%\cite{Fodor:2012td}
\bibitem{Fodor:2012td} 
  Z.~Fodor, K.~Holland, J.~Kuti, D.~Nogradi, and C.~H.~Wong,
  %``The Yang-Mills gradient flow in finite volume,''
  JHEP {\bf 1211}, 007 (2012)
  [arXiv:1208.1051 [hep-lat]].
  %%CITATION = ARXIV:1208.1051;%%
  %21 citations counted in INSPIRE as of 07 Mar 2014

%\cite{Fritzsch:2013je}
\bibitem{Fritzsch:2013je} 
  P.~Fritzsch and A.~Ramos,
  %``The gradient flow coupling in the Schrödinger Functional,''
  JHEP {\bf 1310}, 008 (2013)
  [arXiv:1301.4388 [hep-lat]].
  %%CITATION = ARXIV:1301.4388;%%
  %17 citations counted in INSPIRE as of 07 Mar 2014

%\cite{Monahan:2013lwa}
\bibitem{Monahan:2013lwa} 
  C.~Monahan and K.~Orginos,
  %``Finite volume renormalization scheme for fermionic operators,''
  PoS Lattice {\bf 2013}, 443 (2014)
  [arXiv:1311.2310 [hep-lat]].
  %%CITATION = ARXIV:1311.2310;%%
  %8 citations counted in INSPIRE as of 15 juin 2015

%\cite{Shindler:2013bia}
\bibitem{Shindler:2013bia} 
  A.~Shindler,
  %``Chiral Ward identities, automatic O(a) improvement and the gradient flow,''
  Nucl.\ Phys.\ B {\bf 881}, 71 (2014)
  [arXiv:1312.4908 [hep-lat]].
  %%CITATION = ARXIV:1312.4908;%%
  %1 citations counted in INSPIRE as of 25 Jun 2014

%\cite{Bonati:2014tqa}
\bibitem{Bonati:2014tqa} 
  C.~Bonati and M.~D'Elia,
  %``Comparison of the gradient flow with cooling in $SU(3)$ pure gauge theory,''
  Phys.\ Rev.\ D {\bf 89}, 105005 (2014)
  [arXiv:1401.2441 [hep-lat]].
  %%CITATION = ARXIV:1401.2441;%%
  %1 citations counted in INSPIRE as of 25 Jun 2014

%\cite{DelDebbio:2013zaa}
\bibitem{DelDebbio:2013zaa} 
  L.~Del Debbio, A.~Patella, and A.~Rago,
  %``Space-time symmetries and the Yang-Mills gradient flow,''
  JHEP {\bf 1311}, 212 (2013)
  [arXiv:1306.1173 [hep-th]].
  %%CITATION = ARXIV:1306.1173;%%
  %4 citations counted in INSPIRE as of 07 Mar 2014

%\cite{Asakawa:2013laa}
\bibitem{Asakawa:2013laa} 
  M.~Asakawa {\it et al.}  [FlowQCD Collaboration],
  %``Thermodynamics of SU(3) gauge theory from gradient flow on the lattice,''
  Phys.\ Rev.\ D {\bf 90}, no. 1, 011501 (2014)
  [arXiv:1312.7492 [hep-lat]].
  %%CITATION = ARXIV:1312.7492;%%
  %20 citations counted in INSPIRE as of 15 juin 2015

%\cite{Collins:1984xc}
\bibitem{Collins:1984xc} 
  J.~C.~Collins,
  {\it Renormalization.\ An Introduction to Renormalization, the Renormalization Group, and the Operator Product Expansion\/}
  (Cambridge University Press, Cambridge, 1984).
  %8 citations counted in INSPIRE as of 07 Mar 2014

%\cite{Crewther:1972kn}
\bibitem{Crewther:1972kn} 
  R.~J.~Crewther,
  %``Nonperturbative evaluation of the anomalies in low-energy theorems,''
  Phys.\ Rev.\ Lett.\  {\bf 28}, 1421 (1972).
  %%CITATION = PRLTA,28,1421;%%
  %412 citations counted in INSPIRE as of 07 Mar 2014

%\cite{Chanowitz:1972vd}
\bibitem{Chanowitz:1972vd} 
  M.~S.~Chanowitz and J.~R.~Ellis,
  %``Canonical Anomalies and Broken Scale Invariance,''
  Phys.\ Lett.\ B {\bf 40}, 397 (1972).
  %%CITATION = PHLTA,B40,397;%%
  %261 citations counted in INSPIRE as of 07 Mar 2014

%\cite{Adler:1976zt}
\bibitem{Adler:1976zt} 
  S.~L.~Adler, J.~C.~Collins and A.~Duncan,
  %``Energy-Momentum-Tensor Trace Anomaly in Spin 1/2 Quantum Electrodynamics,''
  Phys.\ Rev.\ D {\bf 15}, 1712 (1977).
  %%CITATION = PHRVA,D15,1712;%%
  %244 citations counted in INSPIRE as of 07 Mar 2014

%\cite{Nielsen:1977sy}
\bibitem{Nielsen:1977sy} 
  N.~K.~Nielsen,
  %``The Energy Momentum Tensor in a Nonabelian Quark Gluon Theory,''
  Nucl.\ Phys.\ B {\bf 120}, 212 (1977).
  %%CITATION = NUPHA,B120,212;%%
  %200 citations counted in INSPIRE as of 07 Mar 2014

%\cite{Collins:1976yq}
\bibitem{Collins:1976yq} 
  J.~C.~Collins, A.~Duncan, and S.~D.~Joglekar,
  %``Trace and Dilatation Anomalies in Gauge Theories,''
  Phys.\ Rev.\ D {\bf 16}, 438 (1977).
  %%CITATION = PHRVA,D16,438;%%
  %492 citations counted in INSPIRE as of 07 Mar 2014

%\cite{Fujikawa:1980rc}
\bibitem{Fujikawa:1980rc} 
  K.~Fujikawa,
  %``Energy Momentum Tensor in Quantum Field Theory,''
  Phys.\ Rev.\ D {\bf 23}, 2262 (1981).
  %%CITATION = PHRVA,D23,2262;%%
  %183 citations counted in INSPIRE as of 07 Mar 2014

%\cite{Fujikawa:2004cx}
\bibitem{Fujikawa:2004cx} 
  K.~Fujikawa and H.~Suzuki,
  {\it Path Integrals and Quantum Anomalies\/} (Clarendon, Oxford, 2004).
  %2 citations counted in INSPIRE as of 07 Mar 2014

%\cite{Caswell:1974gg}
\bibitem{Caswell:1974gg} 
  W.~E.~Caswell,
  %``Asymptotic Behavior of Nonabelian Gauge Theories to Two Loop Order,''
  Phys.\ Rev.\ Lett.\  {\bf 33}, 244 (1974).
  %%CITATION = PRLTA,33,244;%%
  %676 citations counted in INSPIRE as of 07 Mar 2014

%\cite{Jones:1974mm}
\bibitem{Jones:1974mm} 
  D.~R.~T.~Jones,
  %``Two Loop Diagrams in Yang-Mills Theory,''
  Nucl.\ Phys.\ B {\bf 75}, 531 (1974).
  %%CITATION = NUPHA,B75,531;%%
  %592 citations counted in INSPIRE as of 07 Mar 2014

%\cite{Tarrach:1980up}
\bibitem{Tarrach:1980up} 
  R.~Tarrach,
  %``The Pole Mass in Perturbative QCD,''
  Nucl.\ Phys.\ B {\bf 183}, 384 (1981).
  %%CITATION = NUPHA,B183,384;%%
  %299 citations counted in INSPIRE as of 07 Mar 2014

%\cite{Nachtmann:1981zg}
\bibitem{Nachtmann:1981zg} 
  O.~Nachtmann and W.~Wetzel,
  %``The Beta Function for Effective Quark Masses to Two Loops in {QCD},''
  Nucl.\ Phys.\ B {\bf 187}, 333 (1981).
  %%CITATION = NUPHA,B187,333;%%
  %43 citations counted in INSPIRE as of 07 Mar 2014

%\cite{Nakamura:2004sy}
\bibitem{Nakamura:2004sy} 
  A.~Nakamura and S.~Sakai,
  %``Transport coefficients of gluon plasma,''
  Phys.\ Rev.\ Lett.\  {\bf 94}, 072305 (2005)
  [hep-lat/0406009].
  %%CITATION = HEP-LAT/0406009;%%
  %169 citations counted in INSPIRE as of 12 Mar 2014

%\cite{Meyer:2007ic}
\bibitem{Meyer:2007ic} 
  H.~B.~Meyer,
  %``A Calculation of the shear viscosity in SU(3) gluodynamics,''
  Phys.\ Rev.\ D {\bf 76}, 101701 (2007)
  [arXiv:0704.1801 [hep-lat]].
  %%CITATION = ARXIV:0704.1801;%%
  %219 citations counted in INSPIRE as of 12 Mar 2014

%\cite{Meyer:2007dy}
\bibitem{Meyer:2007dy} 
  H.~B.~Meyer,
  %``A Calculation of the bulk viscosity in SU(3) gluodynamics,''
  Phys.\ Rev.\ Lett.\  {\bf 100}, 162001 (2008)
  [arXiv:0710.3717 [hep-lat]].
  %%CITATION = ARXIV:0710.3717;%%
  %148 citations counted in INSPIRE as of 12 Mar 2014

%\cite{Capitani:1998mq}
\bibitem{Capitani:1998mq} 
  S.~Capitani {\it et al.}  [ALPHA Collaboration],
  %``Nonperturbative quark mass renormalization in quenched lattice QCD,''
  Nucl.\ Phys.\ B {\bf 544}, 669 (1999)
  [hep-lat/9810063].
  %%CITATION = HEP-LAT/9810063;%%
  %326 citations counted in INSPIRE as of 07 Mar 2014

%\cite{Itou:2013faa}
\bibitem{Itou:2013faa} 
  E.~Itou,
  %``The twisted Polyakov loop coupling and the search for an IR fixed point,''
  PoS LATTICE {\bf 2013}, 005 (2014)
  [arXiv:1311.2676 [hep-lat]].
  %%CITATION = ARXIV:1311.2676;%%
  %8 citations counted in INSPIRE as of 15 juin 2015

%\cite{Fodor:2014pqa}
\bibitem{Fodor:2014pqa} 
  Z.~Fodor, K.~Holland, J.~Kuti, D.~Nogradi and C.~H.~Wong,
  %``Can a light Higgs impostor hide in composite gauge models?,''
  PoS LATTICE {\bf 2013}, 062 (2014)
  [arXiv:1401.2176 [hep-lat]].
  %%CITATION = ARXIV:1401.2176;%%
  %24 citations counted in INSPIRE as of 15 juin 2015

\end{thebibliography}
